\section{Introduction}
\label{sec:introduction}

An independent set in a graph is a set of pairwise nonadjacent vertices.
For a finite simple graph $G$, write $\alpha(G)$ for the largest size of
an independent set. If $G$ has $n>0$ vertices, its average degree is
$d(G)=2|E(G)|/n$. The graph is $K_r$-free if no $r$ vertices are pairwise
adjacent; throughout, exclusion means exclusion as an ordinary subgraph.
All logarithms are natural. For general graphs, the elementary bound
$\alpha(G)\ge n/(d+1)$ is best possible, as disjoint unions of cliques
show. Excluding a fixed clique forces larger independent sets. We prove
that the improvement is a factor of order $\log d$.

\begin{theorem}\label{thm:main}
For every integer $r\ge4$, there is a constant $c_r>0$ such that every
finite simple $K_r$-free graph $G$ with $n$ vertices and average degree $d\ge2$
satisfies
\[
 \alpha(G)\ge c_r\frac{n\log d}{d}.
\]
\end{theorem}

The constant is uniform in both the number of vertices and the average
degree. The theorem gives a positive answer to the fixed-clique-size
independence conjecture of Ajtai, Erd\H{o}s, Koml\'os and Szemer\'edi,
also recorded as Erd\H{o}s Problem~802. The order $n\log d/d$ is best
possible up to a constant factor even for triangle-free graphs, and
hence for the larger class of $K_r$-free graphs
\cite[pp.~314--315]{AEKS1981}. We do not optimize the constant $c_r$.

\subsection{History and related results}

Ajtai, Koml\'os and Szemer\'edi~\cite{AKS1980} proved the logarithmic
improvement when triangles are excluded.
Shearer~\cite{Shearer1983} subsequently obtained the asymptotic
coefficient~$1$: triangle-free graphs of average degree $d\to\infty$
have independent sets of size at least $(1-o(1))n\log d/d$.
In 1981, Ajtai, Erd\H os, Koml\'os and Szemer\'edi conjectured the same
order of growth for each fixed forbidden clique~\cite[p.~314, Equation~(3)]{AEKS1981}.
They proved the lower bound $\Omega_r(n\log\log d/d)$ as $d\to\infty$,
so excluding any fixed clique already improves on the general bound.
Shearer~\cite[Corollary~2]{Shearer1995} improved this to
$\Omega_r(n\log d/(d\log\log d))$, again in terms of the average
degree. Theorem~\ref{thm:main} removes this remaining $\log\log d$
denominator.

Subsequent refinements describe more of the graph's local structure.
Dutta, Mubayi and Subramanian~\cite{DMS2012} give bounds involving
the full degree sequence. Davies, Kang, Pirot and
Sereni~\cite[Theorem~30 and Corollary~31]{DKPS2020} develop a local-occupancy
framework that improves constants and also gives coloring estimates.
Dhawan~\cite[Theorems~1.3 and~1.4]{DhawanLocallySparse2026} allows bounded
numbers of cliques, either globally or inside neighborhoods.
For a fixed forbidden clique, these estimates retain a $\log\log$
loss in the relevant degree parameter; they provide information beyond
the uniform order sought in Theorem~\ref{thm:main}.

Several stronger local assumptions yield logarithmic bounds.
Alon~\cite[Theorem~1.1 and the remark following its proof]{Alon1996}
treated graphs whose neighborhoods have
bounded chromatic number, obtaining an average-degree bound of the
required order. His argument also admits bounded fractional chromatic
number in each neighborhood. Davies~\cite[Theorems~1 and~3]{Davies2026}
establishes weighted local versions: he bounds individual vertex
inclusion probabilities in a random independent set under bounded
neighborhood maximum average degree or fractional chromatic number.
For every fixed three-colorable forbidden graph, Dhawan, Janzer and
Methuku~\cite[Theorems~1.4 and~1.6]{DJM2025} prove
$(1-o(1))n\log d/d$ as $d\to\infty$, uniformly in $n$.
This attains leading coefficient~$1$ in that restricted class.
Neither bounded neighborhood fractional chromatic number nor exclusion of a fixed
three-colorable graph covers all ordinary $K_r$-free graphs for $r\ge4$.
For the latter restriction, complete tripartite graphs illustrate the
difference: they are $K_4$-free and contain every fixed three-colorable
graph once their parts are large enough.

For comparison, let $r(s,t)$ be the least integer $N$ such that every
$N$-vertex graph contains a $K_s$ or an independent set of size $t$.
The fixed off-diagonal Ramsey estimates in
\cite[Theorem~1.1]{RamseyFive} and \cite[Theorem~1.1]{RamseyFixed}
determine logarithmic exponents for $r(5,t)$ and, respectively,
for $r(s,t)$ with each fixed $s\ge6$.
Together they state
\[
 r(s,t)=\frac{t^{s-1}}{(\log t)^{s-2+o(1)}}
 \qquad\text{for each fixed }s\ge5.
\]
These are estimates for extremal graph orders as $t\to\infty$.
Theorem~\ref{thm:main} controls every input graph using its own average
degree, including graphs whose degree is small relative to their order.
Neither companion is an input to the proof. In particular, the weighted
triangle theorem used here is proved from its own cross-mass hypothesis.


\subsection{The proof in outline}

The proof uses positive vertex weights. For such a vector $w$, write
$w(S)=\sum_{v\in S}w_v$, and define its edge and triangle masses by
\[
 M_G(w)=\sum_{\{u,v\}\in E(G)}w_uw_v,
 \qquad
 T_G(w)=\sum_{\{u,v,z\}\text{ a triangle}}w_uw_vw_z.
\]
Both sums are over unordered sets. The central step is to show that
$T_G(w)$ is at most a constant depending only on $r$ times $M_G(w)$
when $w$ maximizes, over nonnegative vertex weights,
\[
 F_G(w)=\sum_v w_v(1-\log w_v)-M_G(w).
\]
The zero-coordinate convention is $0\log0=0$.
Three stages establish the theorem.

\begin{enumerate}
\item \emph{Extract a bound that survives changes to the graph.}
Section~\ref{sec:optimizer} chooses a maximizing weight and uses random
multipliers of mean one to control the directed edge mass between any
two vertex sets. The multiplier construction is recursive on the excluded
clique size. Section~\ref{sec:triangle} shows that this cross-mass
bound survives edge deletion and splitting a vertex into copies carrying
its original weight, provided each original edge is retained at most once.

\item \emph{Reduce weighted neighborhoods without losing many triangles.}
Sections~\ref{sec:entropy} and~\ref{sec:splitting} prove a triangle bound
for every weighted graph with that cross-mass property. At each reduction
step, we first delete edges whose common neighborhoods have small weight,
charging the lost triangle mass to their edge mass. For each vertex
$u$, walks inside $N(u)$ assign a probability distribution on $N(u)$ to
each starting vertex $v\in N(u)$. The cross-mass bound limits the average
entropy growth of these walks. At a suitable common time, the distributions
started at $v$ and $z$ are close on average over ordered triangles $(u,v,z)$,
weighted by $w_uw_vw_z$.

Sampling these distributions together groups the incident edges at $u$;
close distributions usually give the same group to the two edges of a
triangle. A cutoff on reversed transition probabilities bounds each
group's total neighbor weight; the entropy and density bounds, together
with reversibility, control the average triangle mass discarded by the
cutoff. Splitting vertices according to
the resulting groups reduces the largest neighborhood weight from
$e^x$ to $e^{\sqrt{x}}$. The losses are summable through all iterations.
Once every neighborhood has bounded weight, triangle mass is directly
bounded by edge mass.

\item \emph{Choose independent vertices at controlled cost.}
For a graph with maximum degree at most $\Delta\ge3$, evaluating $F_G$ at
the constant weight $(\log\Delta)/\Delta$ gives an extremal value at
least a constant multiple of $n(\log\Delta)^2/\Delta$.
When we delete a closed neighborhood and retain the old maximizing weights
elsewhere, edges between the chosen vertex's neighbors contribute to
the cost. Averaging over the chosen vertex with probabilities proportional
to its weight counts these edges as triangles.
Section~\ref{sec:closure} uses the triangle bound to find a vertex whose
closed-neighborhood deletion reduces this value by at most
$D_r\log\Delta$. Selecting such a vertex and repeating on the remaining
graph gives an independent set of size at least a constant multiple of
$n\log\Delta/\Delta$. The same degree bound $\Delta$ is kept throughout
this induction. Deleting the original graph's high-degree vertices then
yields the average-degree bound in Theorem~\ref{thm:main}.
\end{enumerate}

The weighted triangle theorem is a separate statement about cross masses;
it has no clique hypothesis. Preservation of the cross-mass bound under
vertex splitting permits repeated reduction without requiring the new weights to
remain maximizers. The entropy estimate controls only averages over a
specified triangle distribution, and the splitting argument uses exactly
those averages. Every threshold in these two steps depends only on the
cross-mass constant. The proof gives an existence result with an
$r$-dependent constant; it makes no claim of an efficient algorithm or
an optimal leading constant.
